\section{Antecedentes y marco teórico}\label{sec:marco}

\subsection{Detección supervisada de pérdidas no técnicas}

La detección de fraude eléctrico con aprendizaje automático tiene unos quince años de
desarrollo. \textcite{nagi2010} sustituyeron la selección de clientes por reglas manuales con
un clasificador de máquinas de vectores de soporte entrenado sobre perfiles de consumo, en
una distribuidora de Malasia. \textcite{comapuig2016} trabajaron con boosting de gradiente
sobre clientes de una empresa española y evaluaron el modelo por campañas de inspección, es
decir, por cuántos fraudes encuentra cada lista que se entrega a los técnicos. En Ecuador,
\textcite{ariasmarin2025} compararon varios clasificadores sobre los datos de una
distribuidora y obtuvieron el mejor resultado con boosting (CatBoost, exactitud de 0,897 y
F1 de 0,894), aunque midieron el desempeño sobre una base de 27.374 registros balanceada
artificialmente entre casos con y sin fraude.

Sobre los mismos datos de STEG que usa este trabajo, \textcite{oprea2022} construyeron
cerca de doscientas variables con funciones analíticas de ventana sobre el historial de
facturas y redujeron el conjunto a quince por puntaje de Fisher. En su escenario de
referencia el mejor AUC fue 0,7331; con la ingeniería de variables propuesta reportan un
AUC de 0,997 para un árbol de decisión. El salto no tiene explicación en el texto, las
variables de ventana se calculan sobre el conjunto completo antes de separar entrenamiento y
prueba, la evaluación descansa en un único corte aleatorio y las quince variables finales no
se nombran. Esas cifras no pueden tomarse como referencia de lo alcanzable. La propia
competencia de Zindi, evaluada con AUC, tuvo un liderato público cercano a 0,86, bajo un
protocolo distinto y sin documentación metodológica publicada por los participantes.

\subsection{La etiqueta depende de la inspección}

En estos conjuntos de datos la etiqueta positiva proviene de una inspección que confirmó el
fraude. \textcite{comapuig2022} hacen de esa procedencia el eje de su trabajo con los datos
de Naturgy: las campañas se dirigen a quien ya parece sospechoso, por lo que el conjunto
etiquetado no representa a la población que se va a puntuar. Documentan un clasificador cuyas
ocho variables más influyentes eran casi todas historial de visitas, es decir, un modelo que
había aprendido el proceso de inspección. \textcite{glauner2016} señalan el mismo sesgo y su
interacción con el desbalance de clases.

De ahí se derivan dos mecanismos que este trabajo mide en sus propios datos. El primero es
la exposición: un cliente con más historial ha tenido más oportunidades de ser
inspeccionado y, por tanto, de ser etiquetado como fraude. El segundo es la huella: si la
inspección confirma el fraude, la empresa reemplaza el contador o modifica el contrato, y esos
cambios quedan en el historial que luego se usa para construir variables.

\subsection{Fuga de información}

Hay fuga cuando el modelo se entrena con información que no estaría disponible en el momento
de predecir \parencite{kaufman2012}. En este problema aparece de dos formas. Una es
procedimental: calcular estadísticos que comparan a un cliente con otros, o escalar y
codificar, usando también los clientes de prueba antes de partir, que es el patrón observado
en \textcite{oprea2022}. La otra es temporal. Como la etiqueta no tiene fecha, cualquier
variable calculada sobre el historial completo puede incluir lo ocurrido después de la
detección, que no existe cuando se puntúa a un cliente todavía no inspeccionado.

\subsection{Métricas para una lista de inspección}

Con una clase minoritaria del 5,58\,\%, un clasificador que nunca predice fraude acierta el
94,42\,\% de las veces, por lo que la exactitud no informa nada. El ROC-AUC mide la
separación global entre clases y no depende de la prevalencia, lo que lo hace adecuado para
comparar subgrupos con tasas de fraude distintas. El PR-AUC, calculado como precisión media,
concentra su peso en la parte alta del ordenamiento, que es donde opera una lista de
inspección, y su valor de referencia al azar es la prevalencia \parencite{saito2015}.

Para una lista con los $k$ clientes de mayor puntuación, con $P$ positivos en la población
y $N$ clientes en total, se usan
\begin{equation}
  \text{precisión@}k = \frac{\text{TP}_k}{k}, \qquad
  \text{recall@}k = \frac{\text{TP}_k}{P}, \qquad
  \text{lift@}k = \frac{\text{precisión@}k}{P/N},
\end{equation}
donde $\text{TP}_k$ es el número de fraudes dentro de la lista. Una lista de tamaño $k < P$
no puede recuperar todos los fraudes, así que el recall@$k$ se compara con su techo
$\min(k, P)/P$ y no con la unidad.

\subsection{Modelos e interpretabilidad}

Los modelos de boosting de gradiente sobre árboles combinan de forma aditiva árboles poco
profundos, cada uno ajustado al error de los anteriores. LightGBM \parencite{ke2017} y XGBoost
\parencite{chen2016} son sus implementaciones más usadas en datos tabulares; ambas toleran
valores faltantes sin imputar y admiten variables categóricas de forma nativa, dos
propiedades relevantes cuando muchos clientes no tienen contador de gas o tienen historiales
cortos. La regresión logística con regularización L2 sirve como referencia lineal.

Para interpretar el modelo se usan los valores SHAP, que reparten la predicción de cada
cliente entre sus variables de forma aditiva \parencite{lundberg2017}; en modelos de árboles
pueden calcularse de manera exacta \parencite{lundberg2020}. Los intervalos de confianza de
las métricas se obtienen por bootstrap \parencite{efron1979} remuestreando clientes, y los de
proporciones del análisis exploratorio con el método de \textcite{wilson1927}.

\subsection{Vacío que aborda este trabajo}

Las técnicas de clasificación avanzaron desde las máquinas de vectores de soporte hasta el
boosting, pero los trabajos revisados rara vez evalúan con la distribución real de clases,
no verifican si la etiqueta depende de la inspección y, en el único antecedente publicado
sobre estos datos, no controlan la fuga. Ninguno separa la parte del desempeño que procede
del comportamiento del cliente de la que procede del rastro de la inspección. Este trabajo
mide esa separación sobre los datos de STEG y evalúa el modelo como se usaría: en listas de
tamaño fijo y frente a referencias simples.
